\chapter{The Bidirectional Equivalence Theorem}
\label{ch:bidirectional-equivalence}

\section{Overview}

Parts~I and II have developed two complete semantic universes:

\begin{itemize}
\item the continuous RSVP universe of scalar--vector--entropy fields
  $(\Phi,\mathbf{v},S)$ evolving on a stratified manifold $\mathcal{M}$ through
  curvature-matched harmonic flows, and

\item the discrete Polyxan universe of typed hypergraphs $\mathcal{G}$ evolving
  through curvature-decreasing EHR rewriting and stabilising as media quines.
\end{itemize}

The present chapter proves the central structural result of Part~II:  

\begin{quote}
\textbf{The RSVP and Polyxan universes are equivalent.}
\end{quote}

More precisely, we construct a pair of functors
\[
\mathbb{R} : \mathbf{RSVP}^{\mathrm{harm}} \longrightarrow \mathbf{Polyx}^{\mathrm{quine}},
\qquad
\mathbb{L} : \mathbf{Polyx}^{\mathrm{quine}} \longrightarrow 
\mathbf{RSVP}^{\mathrm{harm}},
\]
and show that:

\begin{enumerate}
\item they are \emph{quasi-inverse symmetric monoidal $\infty$-equivalences},
\item they are \emph{dynamically faithful}: they intertwine RSVP gradient flow
  and Polyxan rewriting,
\item they \emph{preserve modal structure}: harmonicity corresponds to
  quine-stability,
\item they \emph{preserve curvature stratification}: quine depth corresponds to
  geometric curvature depth,
\item they \emph{share} the unique universal attractor: the universal media
  quine~$\mathfrak{U}$ corresponds under both functors to the conformal vacuum.
\end{enumerate}

This equivalence is the structural basis for the synthetic duality of Part~IV.



\section{The Categories \texorpdfstring{$\mathbf{RSVP}^{\mathrm{harm}}$}{RSVP\^harm} and \texorpdfstring{$\mathbf{Polyx}^{\mathrm{quine}}$}{Polyx\^quine}}

\subsection{Objects}

\paragraph{Harmonic RSVP objects.}
An object of $\mathbf{RSVP}^{\mathrm{harm}}$ is a triple
\[
(\mathcal{M}, (\Phi,\mathbf{v},S), \mathcal{R})
\]
satisfying:

\begin{itemize}
\item $\mathcal{M}$ a stratified RSVP manifold;
\item $(\Phi,\mathbf{v},S)$ a harmonic solution of the RSVP field equations;
\item $\mathcal{R}$ the resulting scalar curvature, integer-valued on strata.
\end{itemize}

\paragraph{Quine-stable Polyxan objects.}
An object of $\mathbf{Polyx}^{\mathrm{quine}}$ is a typed hypergraph
\[
\mathcal{G} = (V,E,\mathscr{T},\mathcal{H}_0)
\]
that is a fixed point of the quine monad:
\[
\mathcal{Q}(\mathcal{G}) \simeq \mathcal{G}.
\]

\subsection{Morphisms}

\paragraph{RSVP morphisms.}
A morphism in $\mathbf{RSVP}^{\mathrm{harm}}$ is a stratified
curvature-preserving smooth map
\[
f : (\mathcal{M}_1,(\Phi_1,\mathbf{v}_1,S_1)) \to
    (\mathcal{M}_2,(\Phi_2,\mathbf{v}_2,S_2))
\]
such that $f^\ast \mathcal{R}_2 = \mathcal{R}_1$.

\paragraph{Polyxan morphisms.}
A morphism in $\mathbf{Polyx}^{\mathrm{quine}}$ is a structure-preserving map
of typed hypergraphs
\[
h : \mathcal{G}_1 \to \mathcal{G}_2
\]
that commutes with the quine structure.



\section{Definition of the Functors \texorpdfstring{$\mathbb{R}$}{R} and \texorpdfstring{$\mathbb{L}$}{L}}

\subsection{The Functor \texorpdfstring{$\mathbb{R}$}{R}: RSVP $\to$ Polyx}

Given a harmonic RSVP configuration
\[
(\mathcal{M},\Phi,\mathbf{v},S),
\]
we define $\mathbb{R}(\mathcal{M})$ to be the Polyxan hypergraph obtained by:

\begin{enumerate}
\item taking the \emph{stratified Morse complex} of $\mathcal{M}$,
\item assigning to each cell a type from $\mathscr{T}$ determined by curvature
  and entropy invariants,
\item defining hyperedges from gradient-flow incidence relations,
\item assigning primitive curvature $\mathcal{H}_0$ from stratum index.
\end{enumerate}

On morphisms, $\mathbb{R}(f)$ is induced functorially on incidence relations.

\subsection{The Functor \texorpdfstring{$\mathbb{L}$}{L}: Polyx $\to$ RSVP}

Given a quine-stable Polyxan hypergraph $\mathcal{G}$,
\[
\mathbb{L}(\mathcal{G})
\]
is defined by the \emph{canonical harmonic embedding}:

\begin{enumerate}
\item take the incidence complex $\mathrm{Inc}(\mathcal{G})$,
\item embed it into an RSVP manifold using the curvature-matching embedding
  theorem (Chapter~31),
\item evolve under RSVP flow to convergence,
\item obtain a harmonic configuration $(\mathcal{M},\Phi,\mathbf{v},S)$ that is
  unique up to diffeomorphism.
\end{enumerate}



\section{Statement of the Bidirectional Equivalence Theorem}

\begin{theorem}[Bidirectional Equivalence]
\label{thm:bidirectional-equivalence}
The functors $\mathbb{R}$ and $\mathbb{L}$ form a quasi-inverse equivalence of
symmetric monoidal $\infty$-categories:
\[
\mathbb{L} \circ \mathbb{R} \;\simeq\; \mathrm{id}_{\mathbf{RSVP}^{\mathrm{harm}}},
\qquad
\mathbb{R} \circ \mathbb{L} \;\simeq\; \mathrm{id}_{\mathbf{Polyx}^{\mathrm{quine}}}.
\]
Moreover, the equivalence is:
\begin{enumerate}
\item \emph{dynamically faithful}:
  RSVP gradient flow corresponds to EHR rewriting,
\item \emph{modally faithful}:
  harmonicity corresponds to quine-stability,
\item \emph{curvature-faithful}:
  curvature strata correspond to quine strata,
\item \emph{attractor-faithful}:
  the universal attractor $(\mathfrak{U},\mathrm{vac})$ is fixed by both.
\end{enumerate}
\end{theorem}



\section{Proof of Categorical Equivalence}

\subsection{Naturality}

We construct natural isomorphisms
\[
\eta : \mathrm{id} \Rightarrow \mathbb{L}\circ\mathbb{R},
\qquad
\epsilon : \mathbb{R}\circ\mathbb{L} \Rightarrow \mathrm{id},
\]
by:

\begin{itemize}
\item $\eta$ mapping each RSVP configuration to the harmonic configuration
  obtained from its own stratified Morse complex,
\item $\epsilon$ mapping each Polyxan hypergraph to the hypergraph obtained
  from its own canonical embedding.
\end{itemize}

Both are isomorphisms because stratified Morse complexes and canonical embedding
flows are unique up to canonical choice (Chapter~31).

\subsection{Monoidal Structure}

Both functors preserve disjoint unions and tensor products:
\[
\mathbb{R}(\mathcal{M}_1 \sqcup \mathcal{M}_2)
 = \mathbb{R}(\mathcal{M}_1) \otimes \mathbb{R}(\mathcal{M}_2),
\]
\[
\mathbb{L}(\mathcal{G}_1 \otimes \mathcal{G}_2)
 = \mathbb{L}(\mathcal{G}_1) \sqcup \mathbb{L}(\mathcal{G}_2).
\]

\subsection{Essential Surjectivity}

Given any harmonic RSVP object
\[
(\mathcal{M},\Phi,\mathbf{v},S),
\]
its Morse complex is a quine-stable Polyxan hypergraph by the curvature-matching
theorem.  
Thus it is in the image of $\mathbb{R}$.

Given any quine-stable Polyxan hypergraph $\mathcal{G}$, the canonical embedding
and harmonic limit exist (Chapters~31--32), so $\mathbb{L}(\mathcal{G})$ is
harmonic.  
Thus it is in the image of $\mathbb{L}$.



\section{Dynamical Equivalence}

\subsection{RSVP Flow Corresponds to EHR Rewriting}

The embedding flow equation (Chapter~32) shows that:

\[
\partial_t \iota
= -\frac{\delta \mathcal{E}}{\delta \iota}
\]
produces topological events exactly matching the EHR rewrite rules.  
The morphogenesis theorem (Chapter~33) proves the bijective correspondence.

Thus the diagram commutes:

\[
\begin{tikzcd}
\mathbf{RSVP}^{\mathrm{flow}} 
  \arrow[r,"\mathbb{R}"]
  \arrow[d,"\mathrm{flow}"']
&
\mathbf{Polyx}^{\mathrm{EHR}}
  \arrow[d,"\mathrm{rewrite}"]
\\
\mathbf{RSVP}^{\mathrm{flow}}
  \arrow[r,"\mathbb{R}"']
&
\mathbf{Polyx}^{\mathrm{EHR}}
\end{tikzcd}
\]

\subsection{Quine Tower ↔ Harmonic Descent}

The quine tower stabilises at height three (Chapter~25).  
The harmonic descent stabilises in finite time (Chapter~32).  
Under $\mathbb{R}$ and $\mathbb{L}$, both processes correspond to each other.



\section{Modal Equivalence}

\begin{theorem}[Modal Correspondence]
For any proposition $P$ about RSVP or Polyx objects,
\[
\Box_{\mathrm{RSVP}} P 
\quad\Longleftrightarrow\quad
\Box_{\mathrm{Polyx}} P
\]
under the translations $\mathbb{R}$ and $\mathbb{L}$.
\end{theorem}

The proof follows from the definitions of modal operators and functorial
preservation of harmonicity and quine-stability.



\section{Universal Attractor Equivalence}

\begin{theorem}[Universal Attractor Correspondence]
The universal media quine $\mathfrak{U}$ corresponds under both functors to the
RSVP conformal vacuum:
\[
\mathbb{L}(\mathfrak{U}) = \mathrm{vac},
\qquad
\mathbb{R}(\mathrm{vac}) = \mathfrak{U}.
\]
\end{theorem}

Thus the two universes share the same unique ground state.



\section{Final Cadence}

There are not two semantic universes.  
There is one.

The continuous RSVP manifold and the discrete Polyxan hypergraph are merely two
presentations of the same stratified geometric object; their flows are the same;
their modal invariances are the same; their universal attractor is the same.

\begin{quote}
\textbf{Geometry is syntax.  
Syntax is geometry.  
Meaning is their equivalence.}
\end{quote}

